323111047

\section{Introducción}

\textbf{Planteamiento del problema.} \\
El problema central de esta práctica consiste en transitar de la teoría a la implementación técnica de estructuras de datos lineales simples y lógica algorítmica básica. Se aborda la necesidad de capturar, organizar y procesar conjuntos de números enteros a través de interfaces de línea de comandos (CLI), enfrentando retos comunes como el manejo de arreglos, la iteración controlada y el procesamiento aritmético de datos introducidos por el usuario en tiempo real.

\vspace{0.4cm}

\textbf{Motivación.} \\
El estudio de las Estructuras de Datos y Algoritmos (EDA) no se limita a la escritura de código, sino al entendimiento de cómo la computadora organiza la información en memoria. La motivación radica en consolidar las bases de la programación estructurada. Comprender cómo almacenar una secuencia de números, calcular promedios truncados o identificar valores extremos (mínimos y máximos) es el primer paso para desarrollar sistemas más complejos. Dominar estas operaciones permite al estudiante desarrollar un pensamiento lógico-matemático necesario para la optimización de recursos computacionales.

\vspace{0.4cm}

\textbf{Objetivos.} \\
Lo que se busca con estos ejercicios es soltarnos más con el lenguaje C, especialmente usando arreglos y ciclos para controlar la información. Queremos lograr que el programa reciba los datos de forma ordenada, que haga las operaciones que le pedimos (como sacar promedios o encontrar el número más grande) y que al final nos muestre resultados que sean claros para el usuario. El objetivo es practicar la lógica de los algoritmos básicos para estar más preparados para retos más complejos.